\subsection{An arithmetic comparison operator}

The suggestion that primes, a zeta function, and uncertainty might
participate in the same selection mechanism can be made mathematically
precise without identifying any of them with a physical observable.
This distinction matters: an analogy becomes testable only after its
operators, state space, and domains have been specified. The following
construction is an analytic comparison model. It does not extend the
Companion action or supply a lepton-mass prediction.

Fix $\sigma>1$ and use the positive Hilbert space
$\mathcal H=L^2(\mathbb R,dt)$, with dimensionless $t$. The logarithmic
derivative of the Euler product gives \cite{DLMF274}
\begin{equation}
 -\frac{\zeta'(\sigma+it)}{\zeta(\sigma+it)}
 =\sum_{n\geq2}\Lambda(n)n^{-\sigma}e^{-it\log n},
 \qquad
 \Lambda(p^k)=\log p.
\end{equation}
Here $\Lambda(n)=0$ unless $n$ is a prime power. In particular the
real odd function
\begin{equation}
 a(t)=\operatorname{Im}\frac{\zeta'(\sigma+it)}{\zeta(\sigma+it)}
     =\sum_{n\geq2}\Lambda(n)n^{-\sigma}\sin(t\log n)
\end{equation}
defines a bounded self-adjoint multiplication operator $Af=af$.
Absolute and uniform convergence follow from
$\Lambda(n)\leq\log n$ and convergence of
$\sum (\log n)n^{-\sigma}$. Every fixed derivative is also uniformly
convergent; the relevant majorant has an additional fixed power of
$\log n$. Reflection $Rf(t)=f(-t)$ is a self-adjoint involution and
$RA=-AR$. Thus a parity operator and a genuinely arithmetic odd
partner coexist in one positive Hilbert space.

For a normalized vector, Robertson's inequality \cite{Robertson1929}
becomes
\begin{equation}
 \Delta R\,\Delta A\geq
 \tfrac12|\langle[R,A]\rangle|=|\langle RA\rangle|.
\end{equation}
It does not require $\langle R\rangle=0$. An even normalized vector
has $\langle R\rangle=1$ and both $\langle A\rangle$ and
$\langle RA\rangle$ zero, so it satisfies the inequality without
balance. This is a direct counterexample to an uncertainty-only
selection argument. The inequality restricts a prepared state; it
does not choose the preparation.

\subsection{Why an eigenstate prescription fails}

The bounded operator $T=-iRA$ is self-adjoint and satisfies
$T^2=A^2$ and $RT=-TR$. If a normalized eigenvector with nonzero
eigenvalue existed, anticommutation would force its parity expectation
to vanish. That observation suggests a stronger selection rule than
Robertson alone. In this particular realization, however, there are
no nonzero normalizable eigenvectors, including at eigenvalue zero.

Indeed, the sine series is holomorphic in the strip
$|\operatorname{Im}t|<\sigma-1$ and hence real analytic on the real
axis. Moreover,
\begin{equation}
 a(0)=0,\qquad
 a'(0)=\sum_{n\geq2}\Lambda(n)(\log n)n^{-\sigma}>0.
\end{equation}
Consequently $a^2$ is not constant. An eigenvector of $T$ at real
$\lambda$ would obey $(a(t)^2-\lambda^2)\psi(t)=0$ almost
everywhere. The zero set of that nontrivial real-analytic function
has measure zero, and an $L^2$ vector supported on it is zero.
The argument excludes this eigenstate route, rather than every
possible role for arithmetic. A finite grid could display discrete
eigenvectors while missing this obstruction in the continuum.

\subsection{A regular balanced spectral packet}

One can retain the positive continuous spectral subspace instead.
Let $S=\operatorname{sgn}(A)$; the null set of $a$ is negligible,
so $S^2=I$ and $RS=-SR$. Choose the real even window
$g(t)=e^{-t^2}$, set $f=Na^2g$ with $\|f\|=1$, and define
\begin{equation}
 \psi=\frac{f+iSf}{\sqrt2},\qquad
 \operatorname{sgn}(T)=-iRS.
\end{equation}
Since $f$ is even and $Sf$ odd, their norms agree and their inner
product vanishes. Direct substitution gives
\begin{equation}
 \|\psi\|=1,\quad \langle R\rangle_\psi=0,\quad
 \operatorname{sgn}(T)\psi=\psi,\quad T\psi=|a|\psi.
\end{equation}
The last expression is multiplication by a variable function,
not an eigenvalue equation. Thus it is consistent with the preceding
nonexistence result. If $F_+=\mathbf1_{(0,\infty)}(T)$, then
$F_+f=(f+iSf)/2$ and normalization gives this packet.

The factor $a^2$ also resolves a domain issue. A generic sign change
in $Sf$ could create a jump and an infinite Fourier second moment.
Here $Sf=Na|a|g$, with weak derivatives
\begin{align}
 f'&=N(2aa'g+a^2g'),\\
 (Sf)'&=N(2|a|a'g+a|a|g').
\end{align}
They lie in $L^2$ because $a,a'$ are bounded and the Gaussian and its
derivative are square integrable. There is no delta contribution at a
zero of $a$: the profile vanishes there. In addition $t\psi$ and
$t\psi'$ lie in $L^2$. With $Qf=tf$ and the usual self-adjoint
$P=-i\,d/dt$, both second moments are finite, $QP\psi$ and
$PQ\psi$ exist, and $[Q,P]\psi=i\psi$. Hence
\begin{equation}
 \Delta Q\,\Delta P\geq\tfrac12.
\end{equation}
This is a parameter--Fourier uncertainty relation. Physical position,
momentum, and a dimensional value of $\hbar$ have not been identified.

Balance does not mean uncertainty saturation. Write
$\mu_1=\int |a||f|^2dt$ and $\mu_2=\int a^2|f|^2dt$. For this packet
\begin{equation}
 \langle R\rangle=\langle A\rangle=0,\quad
 \Delta R=1,\quad \Delta A=\sqrt{\mu_2},\quad
 \langle RA\rangle=i\mu_1.
\end{equation}
The Robertson bound is $\sqrt{\mu_2}\geq\mu_1$ and is strict:
equality would require $|a|$ constant almost everywhere on the
support of $f$, whereas its level sets have measure zero. Balance
comes from the chosen positive spectral projection. It is not an
effect of saturating that inequality.

\subsection{Preparation, preservation, and physical identification}

This comparison separates three operations that an effective theory
must not conflate. A spectral projector defines an allowed subspace;
a preparation mechanism produces states in it; a generator preserves
it during evolution. Functional calculus ensures that $e^{-isT}$
commutes with $F_+$, but this does not prepare a generic incoming
state in $\operatorname{im}F_+$, nor identify $T$ with the native
Hamiltonian. The window, its scale, and the suppression of zeros
are additional choices. None follows from the Euler product.

The construction is confined to $\sigma>1$. Continuing the scalar
zeta function does not automatically continue this bounded operator
or its domain. On the critical line, poles of the logarithmic
derivative require a separate analysis. No statement here proves
the Riemann hypothesis or turns its zeros into the observed masses.
The useful outcome is narrower: an arithmetic odd partner and a
regular balanced packet are compatible with uncertainty, while
their physical preparation and mass readout remain open.

\subsection{What would constitute progress}

For the mass branch, a decisive next result would specify the native
operator $\widetilde R$, its compressed transition $B$, the positive
three-dimensional readout, and the rule selecting its phase. A fit
to three masses tests the flexibility of a chosen parameterization;
it does not supply these missing definitions. A scale-independent
ratio can constrain a mechanism before an overall mass scale is
known, but agreement of a ratio cannot identify that scale or the
underlying dynamics.

For the generation branch, an admissible seed map must retain rank
three and satisfy the grading requirements through the actual
operator chain. Mixing requires a second physically identified
operator in the same comparison space. For the geometric branch,
an invertible coframe, compatible finite derivative data, and a
local solution are separate obligations. Constraint propagation is
valuable after its field-equation hypotheses are met.

The database edition makes these distinctions inspectable. It
stores statement versions, dependencies, source hashes, and the
editorial selection used to assemble this paper. Its purpose is to
prevent an algebraic possibility, a conditional calculation, and
a physical identification from becoming indistinguishable during
rewriting. The present work supplies explicit constructions and
obstructions for these interfaces. A completed phenomenological
theory would additionally need independently specified observables,
units, preparation rules, and comparison with measured data.

\subsection{Open research question: expansion and unequal string relaxation}

The following comparison records the open question at the v5 release
boundary. It is an analytical exploration, not a native consequence of
the graded carrier or a physical identification of fundamental strings.
Suppose three extensible inner paths have lengths $\ell_i$, while an
outer path has fixed length $L$. The equality $\sum_i\ell_i=L$ can be
imposed throughout motion or only as a terminal condition; these are
different models. A concrete network, its winding classes and its
endpoint constraints have not yet been specified.

\paragraph{Elastic allocation and its limitation.}
For independent nonnegative lengths, minimizing
$U=\tfrac12\sum_i k_i\ell_i^2$, with $k_i>0$ and fixed sum $L$,
gives $\ell_i/L=k_i^{-1}/\sum_j k_j^{-1}$ and equal generalized
tensions. Under the additional readout $\sqrt{m_i}=c\ell_i$, $c>0$,
the Koide ratio is $K=\sum_i\ell_i^2/L^2$.
Equal stiffnesses give $K=1/3$. Stiffnesses $(k/q,k,k)$ give
$K=(q^2+2)/(q+2)^2$, so $q=4+3\sqrt2$ yields $K=2/3$.
This factor is inferred from the target, not independently predicted;
the example also has two equal masses. Moreover, $K>1/2$ requires one
length greater than $L/2$. The paths therefore cannot be the three
sides of a straight triangle or three pairwise shortest geodesics in
one metric. Prescribed winding, guides or another network would be needed.

\paragraph{A length scale associated with a cosmological constant.}
Dimensional consistency permits $\Lambda=C/L^2$, but does not determine
$C$ or a stress--energy tensor. If the outer loop is additionally
identified with a great circle on a pure four-dimensional de Sitter
horizon, its radius $R_\Lambda=\sqrt{3/\Lambda}$ gives
$L=2\pi R_\Lambda$ and $C=12\pi^2$.
This fixed horizon scale is not a material loop comoving with expansion.
If instead $L=Aa(t)$ and a separately conserved homogeneous,
isotropic component has $\rho=B/L^2$, with constants $A,B>0$, the continuity equation
$\dot\rho+3H(\rho+p)=0$ gives $w=p/\rho=-1/3$, in units $c=1$.
Vacuum energy requires $w=-1$. Three individual paths do not by
themselves supply the required isotropic cosmological component.

\paragraph{A dissipative comparison.}
Let $r_i$ be an evolving unloaded length. In the tensile regime, a
Maxwell-type model with constant $k_i,\eta_i>0$ is
\begin{equation}
 F_i=k_i(\ell_i-r_i),\qquad
 \dot r_i=F_i/\eta_i,\qquad
 \dot F_i=k_i\dot\ell_i-F_i/\tau_i,\quad \tau_i=\eta_i/k_i.
\end{equation}
Different relaxation times allow different responses to the same imposed
extension. Tension can decrease while length increases, provided
$k_i\dot\ell_i<F_i/\tau_i$. At fixed length,
$F_i(t)=F_i(0)e^{-t/\tau_i}$. The elastic energy satisfies
\begin{equation}
 \dot U_i=F_i\dot\ell_i-F_i^2/\eta_i,\qquad
 U_i=\tfrac12 k_i(\ell_i-r_i)^2.
\end{equation}
The dissipated energy must enter the full energy accounting. Slack paths
require a unilateral tensile law; negative spring force must not be
interpreted as compressive string tension. This comparison represents
creep and relaxation, not a derived cyclic-fatigue law or microscopic
damage mechanism.

If $\sum_i\ell_i=L$ holds at every instant, then
$\sum_i\dot\ell_i=0$: all three actual lengths cannot grow together,
although their unloaded lengths may grow and their tensions may relax.
For freely redistributed quasistatic lengths, force balance can equalize
tensions even when material responses differ. A soft penalty
$\kappa(L-\sum_i\ell_i)^2/2$ supplies the same generalized force to
each length, and does not impose exact equality of the lengths' sum.

\paragraph{The retained question.}
Can a justified coupling between expansion, internal relaxation and a
fixed outer length select a stable unequal allocation? Answering it
requires an embedded network, an expansion-to-strain map, an energy
balance including dissipative degrees of freedom, and a mechanism that
limits relaxation. Identical initial data and identical deterministic
laws alone do not select distinct responses. No value of $\Lambda$,
lepton spectrum, or dynamically selected $K=2/3$ follows here. All
existing physical-identification guards remain open.

\paragraph{Sources for the comparison.}
The de Sitter radius is discussed by H.-Y.~Guo,
\emph{The Beltrami Model of De Sitter Space} (2006),
arXiv:hep-th/0607017. For vacuum cosmology see S.~M.~Carroll,
\emph{A No-Nonsense Introduction to General Relativity} (2001),
Section~6. For finite-rate viscoelastic comparisons see C.-Y.~Lin
et al., \emph{Polymers} \textbf{14} (2022), 2124,
doi:10.3390/polym14102124. These sources support the background models,
not their proposed identification with the Companion's physical sectors.
